% 问题二：求解过程

\textbf{4 种方法求解结果对比如表~\ref{tab:q2_compare} 所示。}

\begin{table}[htbp]
\centering
\caption{问题 2 四种方法求解结果对比}
\label{tab:q2_compare}
\begin{tabular}{crrrrr}
\toprule
方法 & 撤销数 & A 撤销 & B 撤销 & C 撤销 & 调整数 \\
\midrule
贪心+优先级 & 53 & 0 & 3 & 50 & 61 \\
整数规划 ILP & 71 & 1 & 18 & 52 & 59 \\
贪心+局部搜索 & 53 & 0 & 3 & 50 & 61 \\
模拟退火 SA & \textbf{50} & 0 & 3 & 47 & 65 \\
\bottomrule
\end{tabular}
\end{table}

\textbf{方案 1：贪心 + 优先级}

按 A > B > C 优先级排序，同类内按冲突度数降序。对每个计划依次尝试：保留（与已保留计划 0 冲突）→ 频段平移（$\pm 1$ 到 $\pm 10$）→ 时间平移（$\pm 1$ 到 $\pm 5$）→ 撤销。每档内按幅度从小到大枚举，选择最小幅度的可行动作。

结果：撤销 53 个（A:0, B:3, C:50），调整 61 个，消解后冲突数为 0。

\textbf{方案 2：整数线性规划 ILP}

将每个计划的动作选择建模为 0-1 决策变量，以撤销数最小化为目标函数，用 CBC 求解器求解。由于变量规模大（150 个计划 $\times$ 多种动作），求解器在 120 秒超时内未能收敛到全局最优，得到的是可行解而非最优解。

结果：撤销 71 个（A:1, B:18, C:52），效果最差。ILP 效果差的原因在于动作空间过大导致松弛困难，求解器时间限制内无法找到高质量解。

\textbf{方案 3：贪心 + 局部搜索}

在方案 1 的贪心解基础上，对相邻计划进行动作交换的局部搜索。由于贪心解已较优，局部搜索未能进一步改善。

结果：与方案 1 完全相同（撤销 53，调整 61）。

\textbf{方案 4：模拟退火 SA}

以贪心解为初始解，采用以下策略：
\begin{itemize}
    \item 随机扰动：每次随机挑选 1\textasciitilde3 个计划，为其更换合法动作；
    \item 多种子策略：使用 $\geq 5$ 个不同随机种子各独立运行，取最优结果；
    \item 收尾消毒：退火结束后逐对检测残余冲突，撤销每对中优先级较低的计划，保证最终 0 冲突。
\end{itemize}

结果：撤销 \textbf{50} 个（A:0, B:3, C:47），调整 65 个。撤销数为 4 种方法中最少。

\textbf{最优方案选择}

模拟退火（SA）的撤销数最少（50 个），且 A 类撤销为 0，符合优先级要求。以 SA 方案作为问题 2 的最终输出，消解后统计结果如表~\ref{tab:q2_result} 所示。

\begin{table}[htbp]
\centering
\caption{时频冲突消解的统计结果}
\label{tab:q2_result}
\begin{tabular}{crrr}
\toprule
类别 & 保留数量 & 调整数量 & 撤销数量 \\
\midrule
A 类 & 19 & 1 & 0 \\
B 类 & 7 & 30 & 3 \\
C 类 & 9 & 34 & 47 \\
\bottomrule
\end{tabular}
\end{table}
